% ============================ WORKSHEET 3 ==================================
\worksheetheading{3}{Quantum numbers and atomic orbitals}
\data{$c = \SI{3.00e8}{\metre\per\second}$\quad $h = \SI{6.62e-34}{\joule\second}$\quad
$1\ \mathrm{eV} = \SI{1.6e-19}{\joule}$\quad $a_0 = 52.9$~pm\quad
hydrogen-like species: $E_n = -13.6\,Z^2/n^2$~eV, $r = n^2a_0/Z$. Calculator allowed.}

\wq{1}{The third shell}
For $n = 3$, list the allowed values of $\ell$ with the name of each subshell, and the values of
$m_\ell$ for each. How many orbitals does the shell contain?
\begin{solution}
$\ell = 0$ (3s): $m_\ell = 0$.\quad $\ell = 1$ (3p): $m_\ell = -1, 0, +1$.\quad
$\ell = 2$ (3d): $m_\ell = -2, -1, 0, +1, +2$.\quad Total $1 + 3 + 5 = 9 = 3^2$ orbitals.
\end{solution}

\wq{1}{Allowed or not?}
Which of these sets $(n, \ell, m_\ell)$ are allowed? Name the orbital when it is, and give the reason
when it is not: (a) $(2, 2, 0)$; (b) $(3, 1, -1)$; (c) $(1, 0, 1)$; (d) $(4, 3, -2)$;
(e) $(2, 1, 2)$; (f) $(5, 0, 0)$.
\begin{solution}
(a) No: $\ell\le n-1 = 1$.\quad (b) Yes, 3p.\quad (c) No: $\ell = 0$ forces $m_\ell = 0$.\quad
(d) Yes, 4f.\quad (e) No: $|m_\ell|\le\ell = 1$.\quad (f) Yes, 5s.
\end{solution}

\wq{1}{Naming wavefunctions}
Name the orbital described by $\Psi_{2,0,0}$, $\Psi_{3,1,1}$, $\Psi_{4,2,-2}$ and $\Psi_{4,3,0}$.
Conversely, give all the wavefunctions of the 2p subshell.
\begin{solution}
2s, 3p, 4d, 4f.\quad 2p: $\Psi_{2,1,-1}$, $\Psi_{2,1,0}$, $\Psi_{2,1,1}$.
\end{solution}

\wq{1}{Spin and boundary surfaces}
(a) Which values can $m_s$ take, and how are they drawn? Which experiment revealed spin?
(b) What is a boundary surface? What is the shape of the boundary surface of an s orbital, and of a p
orbital?
\begin{solution}
(a) $m_s = +\frac12$ ($\uparrow$) or $-\frac12$ ($\downarrow$). The splitting of the sodium D line into
two lines (589.0 and 589.6~nm).\\
(b) A surface enclosing a given probability (usually 95\,\%) of finding the electron. s: a sphere
centred on the nucleus; p: two lobes on either side of the nucleus along one axis (dumbbell), separated by
a nodal plane.
\end{solution}

\wq{2}{Energies and degeneracy in hydrogen}
(a) Give the energy of the electron of hydrogen described by $\Psi_{3,2,-1}$ and by $\Psi_{3,0,0}$.
What word describes this situation? (b) How many degenerate orbitals are there in the shell $n = 4$ of
hydrogen?
\begin{solution}
(a) Both have $n = 3$: $E_3 = -13.6/9 = -1.51$~eV. In hydrogen the energy depends only on $n$: the 3d
and 3s orbitals are \textbf{degenerate}.\\
(b) $n^2 = 16$ (one 4s, three 4p, five 4d, seven 4f).
\end{solution}

\wq{2}{Ground or excited?}
For the hydrogen atom, say which orbital is occupied, whether the atom is in its ground state or an
excited state, and its energy, when the electron is described by (a) $\Psi_{1,0,0}$,
(b) $\Psi_{2,1,0}$, (c) $\Psi_{3,2,1}$. (d) How much energy is needed to go from (a) to (b)?
\begin{solution}
(a) 1s, ground state, $-13.6$~eV.\quad (b) 2p, excited, $-3.40$~eV.\quad (c) 3d, excited, $-1.51$~eV.\\
(d) $-3.40 - (-13.6) = 10.2$~eV (a UV photon of 122~nm).
\end{solution}

\wq{2}{Counting nodes}
Give the number of radial nodes, angular nodes and total nodes of the 2s, 3p, 3d, 4s and 4f orbitals.
Check the 3s, 3p and 3d cases on Figure 3-3 of the handout.
\begin{solution}
Radial $n-\ell-1$, angular $\ell$, total $n-1$:\quad
2s: 1, 0, 1;\quad 3p: 1, 1, 2;\quad 3d: 0, 2, 2;\quad 4s: 3, 0, 3;\quad 4f: 0, 3, 3.\\
On Figure 3-3, the radial curve of 3s falls to zero twice before its main maximum (2 radial nodes),
3p once (1) and 3d never (0).
\end{solution}

\wq{2}{The three p orbitals}
For each of p$_x$, p$_y$ and p$_z$, give the axis along which it points and its nodal plane. Which
one corresponds to $m_\ell = 0$? What happens to the sign of the wavefunction from one lobe to the other?
\begin{solution}
p$_x$: along $x$, nodal plane (yOz).\quad p$_y$: along $y$, nodal plane (xOz).\quad
p$_z$: along $z$, nodal plane (xOy). p$_z$ corresponds to $m_\ell = 0$ (p$_x$ and p$_y$ are combinations of
$m_\ell = \pm1$). The wavefunction changes sign across the nodal plane: the two lobes have opposite
signs (opposite colours).
\end{solution}

\wq{3}{Size of orbitals}
Using $r = n^2a_0/Z$, compute the radius of: (a) the 1s, 2s and 3s orbitals of hydrogen;
(b) the 1s orbital of $\ce{He+}$ and of $\ce{Li^2+}$. (c) Explain in one sentence the effect of $n$ and
of $Z$.
\begin{solution}
(a) H: 1s: 52.9~pm; 2s: $4\times52.9 = 212$~pm; 3s: $9\times52.9 = 476$~pm.\\
(b) $\ce{He+}$ ($Z = 2$): $52.9/2 = 26.5$~pm; $\ce{Li^2+}$ ($Z = 3$): $52.9/3 = 17.6$~pm.\\
(c) The orbital grows as $n^2$ (the electron is less bound and further out) and shrinks as $1/Z$
(a larger nuclear charge pulls the electron closer).
\end{solution}

\wq{3}{Ordering subshells}
Using Madelung's rule, write $n+\ell$ for each subshell and put them in order of increasing energy:
5s, 4f, 3d, 4p, 6s, 4d, 5p, 4s.
\begin{solution}
$n+\ell$: 4s: 4;\ 3d: 5;\ 4p: 5;\ 5s: 5;\ 4d: 6;\ 5p: 6;\ 6s: 6;\ 4f: 7.\\
Order (equal $n+\ell$ sorted by increasing $n$): 4s $<$ 3d $<$ 4p $<$ 5s $<$ 4d $<$ 5p $<$ 6s $<$ 4f.
\end{solution}

\wq{3}{Hydrogen versus potassium}
In hydrogen the 3d orbitals have the same energy as 3s, and lie far below 4s. In potassium
($Z = 19$) the 4s orbital is lower than 3d. (a) Which approximation lets us still use the labels
1s, 2s, 2p \dots\ for many-electron atoms? (b) Explain why the energy now depends on $\ell$ and why 4s
can drop below 3d.
\begin{solution}
(a) The orbital approximation: electron--electron repulsions are neglected in front of the attraction
of the nucleus, so each electron occupies its own hydrogen-like AO.\\
(b) The inner electrons screen the nucleus. An s electron penetrates close to the nucleus (small inner
maxima of its radial density), feels a larger effective charge and is more strongly bound than p and d
electrons of the same shell: $E$ now depends on $\ell$. The 4s orbital penetrates so much that its
energy falls below that of 3d, which stays outside the core; Madelung's rule ($n+\ell = 4$ for 4s,
5 for 3d) summarises this.
\end{solution}

\wq{3}{The Zeeman effect}
(a) Into how many sublevels does a p subshell split in a magnetic field? a d subshell? an f
subshell? (b) Hydrogen atoms emit the $2\mathrm p\to1\mathrm s$ line (122~nm). How many lines are seen in a
magnetic field? (c) What does the Zeeman effect prove?
\begin{solution}
(a) $2\ell+1$ sublevels: p: 3, d: 5, f: 7.\\
(b) 1s does not split ($\ell = 0$); 2p gives 3 sublevels, so 3 lines very close to 122~nm.\\
(c) That the orbitals of a subshell have different orientations, labelled by $m_\ell$ with $2\ell+1$
values; without a field they are degenerate.
\end{solution}

\wq{4}{How many electrons can a shell hold?}
An electron is described by four quantum numbers $(n, \ell, m_\ell, m_s)$. (a) List all the allowed
sets for $n = 2$. (b) Show that the number of sets for a shell $n$ is $2n^2$ and check it for
$n = 1, 2, 3$. (Unit 4: no two electrons of an atom can share the same set, so this is the maximum
number of electrons of the shell.)
\begin{solution}
(a) $(2,0,0,\pm\frac12)$; $(2,1,-1,\pm\frac12)$; $(2,1,0,\pm\frac12)$; $(2,1,1,\pm\frac12)$: 8 sets.\\
(b) Number of orbitals: $\sum_{\ell=0}^{n-1}(2\ell+1) = 2\cdot\frac{(n-1)n}{2} + n = n^2$; each
orbital has two values of $m_s$: $2n^2$. $n = 1$: 2;\ $n = 2$: 8;\ $n = 3$: 18.
\end{solution}

\wq{4}{The sodium doublet}
The D line of sodium is a doublet: $\lambda_1 = 589.0$~nm and $\lambda_2 = 589.6$~nm. (a) Compute the
energy of each photon in J and in eV (4 significant figures). (b) Compute the energy difference
between the two upper levels in J and in meV. (c) Explain the origin of the doublet.
\begin{solution}
(a) $E_1 = hc/\lambda_1 = 6.62\times10^{-34}\times3.00\times10^8/589.0\times10^{-9} = 3.372\times10^{-19}$~J
$= 2.107$~eV;\quad $E_2 = 3.368\times10^{-19}$~J $= 2.105$~eV.\\
(b) $\Delta E = 3.43\times10^{-22}$~J $= 2.1\times10^{-3}$~eV $= 2.1$~meV (about 0.1\,\% of $E$). Keep
all digits until the subtraction: the difference of two close numbers loses precision.\\
(c) The line is the transition 3p$\,\to\,$3s of the outer electron. Its spin ($m_s = \pm\frac12$)
interacts with its orbital motion, which splits 3p into two levels 2.1~meV apart; 3s ($\ell = 0$) is not
split, so there are two lines.
\end{solution}

\wq{4}{Where is the 1s electron?}
The radial probability density of the 1s orbital of hydrogen is
$P(r) = \dfrac{4}{a_0^3}\,r^2\,\mathrm e^{-2r/a_0}$. (a) Show that $P$ is maximum at $r = a_0$.
(b) Explain why the probability density is zero at $r = 0$ even though $|\Psi_{1s}|^2$ is largest
at the nucleus. (c) For $\ce{He+}$, $P(r) \propto r^2\mathrm e^{-2Zr/a_0}$: where is its maximum?
\begin{solution}
(a) $\dfrac{\mathrm dP}{\mathrm dr} = \dfrac4{a_0^3}\,\mathrm e^{-2r/a_0}\left(2r - \dfrac{2r^2}{a_0}\right)
= \dfrac{8r}{a_0^3}\,\mathrm e^{-2r/a_0}\left(1 - \dfrac r{a_0}\right)$, which is zero for $r = a_0$
(positive before, negative after): maximum at $r = a_0 = 52.9$~pm, the Bohr radius.\\
(b) $P(r)$ is the probability per unit distance of being in a thin shell of radius $r$, whose volume
$4\pi r^2\mathrm dr$ tends to zero at the centre. The density per unit volume is largest at the nucleus,
but there is almost no volume there.\\
(c) Same calculation with $a_0/Z$ instead of $a_0$: maximum at $r = a_0/Z = 26.5$~pm, as given by
$r = n^2a_0/Z$.
\end{solution}
